\documentclass[11pt]{article}
% ===================================================================
% Shared preamble for the "tiny lemma, read slowly" preprint series.
% Mirrors the style of FrobeniusLemmas_preprint.tex.
% ===================================================================

\usepackage{amsmath,amssymb,amsthm}
\usepackage{mathtools}
\usepackage[margin=1.2in]{geometry}
\usepackage{hyperref}
\usepackage{xcolor}
\usepackage{listings}
\usepackage{tikz}
\usepackage{tikz-cd}
\usepackage{mathpartir}
\usepackage{newunicodechar}
\usepackage{setspace}

\onehalfspacing

% ---- Unicode glyphs ----
\newunicodechar{→}{\ensuremath{\to}}
\newunicodechar{↦}{\ensuremath{\mapsto}}
\newunicodechar{↔}{\ensuremath{\leftrightarrow}}
\newunicodechar{⟹}{\ensuremath{\Longrightarrow}}
\newunicodechar{⟶}{\ensuremath{\longrightarrow}}
\newunicodechar{≠}{\ensuremath{\neq}}
\newunicodechar{≤}{\ensuremath{\leq}}
\newunicodechar{≥}{\ensuremath{\geq}}
\newunicodechar{∀}{\ensuremath{\forall}}
\newunicodechar{∃}{\ensuremath{\exists}}
\newunicodechar{∘}{\ensuremath{\circ}}
\newunicodechar{·}{\ensuremath{\cdot}}
\newunicodechar{∈}{\ensuremath{\in}}
\newunicodechar{∉}{\ensuremath{\notin}}
\newunicodechar{≃}{\ensuremath{\simeq}}
\newunicodechar{≅}{\ensuremath{\cong}}
\newunicodechar{≈}{\ensuremath{\approx}}
\newunicodechar{∧}{\ensuremath{\wedge}}
\newunicodechar{∨}{\ensuremath{\vee}}
\newunicodechar{¬}{\ensuremath{\neg}}
\newunicodechar{⊢}{\ensuremath{\vdash}}
\newunicodechar{⟨}{\ensuremath{\langle}}
\newunicodechar{⟩}{\ensuremath{\rangle}}
\newunicodechar{⇑}{\ensuremath{\mathord{\uparrow}}}
\newunicodechar{↑}{\ensuremath{\mathord{\uparrow}}}
\newunicodechar{∎}{\ensuremath{\blacksquare}}
\newunicodechar{∣}{\ensuremath{\mid}}
\newunicodechar{∤}{\ensuremath{\nmid}}
\newunicodechar{∑}{\ensuremath{\textstyle\sum}}
\newunicodechar{∏}{\ensuremath{\textstyle\prod}}
\newunicodechar{∅}{\ensuremath{\emptyset}}
\newunicodechar{∪}{\ensuremath{\cup}}
\newunicodechar{∩}{\ensuremath{\cap}}
\newunicodechar{≡}{\ensuremath{\equiv}}
\newunicodechar{ℕ}{\ensuremath{\mathbb{N}}}
\newunicodechar{ℤ}{\ensuremath{\mathbb{Z}}}
\newunicodechar{ℝ}{\ensuremath{\mathbb{R}}}
\newunicodechar{𝔽}{\ensuremath{\mathbb{F}}}
\newunicodechar{α}{\ensuremath{\alpha}}
\newunicodechar{β}{\ensuremath{\beta}}
\newunicodechar{σ}{\ensuremath{\sigma}}
\newunicodechar{τ}{\ensuremath{\tau}}
\newunicodechar{φ}{\ensuremath{\varphi}}
\newunicodechar{ψ}{\ensuremath{\psi}}
\newunicodechar{χ}{\ensuremath{\chi}}
\newunicodechar{Φ}{\ensuremath{\Phi}}
\newunicodechar{Δ}{\ensuremath{\Delta}}
\newunicodechar{₀}{\textsubscript{0}}
\newunicodechar{₁}{\textsubscript{1}}
\newunicodechar{₂}{\textsubscript{2}}
\newunicodechar{ᵏ}{\ensuremath{{}^{k}}}
\newunicodechar{ⁿ}{\ensuremath{{}^{n}}}
\newunicodechar{ʲ}{\ensuremath{{}^{j}}}
\newunicodechar{²}{\ensuremath{{}^{2}}}
\newunicodechar{³}{\ensuremath{{}^{3}}}
\newunicodechar{⁴}{\ensuremath{{}^{4}}}
\newunicodechar{⁻}{\ensuremath{{}^{-}}}
\newunicodechar{¹}{\ensuremath{{}^{1}}}

% ---- Lean listing style ----
\definecolor{leanbg}{RGB}{247,247,250}
\definecolor{leankw}{RGB}{0,90,160}
\definecolor{leancm}{RGB}{120,120,120}
\lstdefinestyle{lean}{
  backgroundcolor=\color{leanbg},
  basicstyle=\ttfamily\small,
  keywordstyle=\color{leankw}\bfseries,
  commentstyle=\color{leancm}\itshape,
  morekeywords={theorem,lemma,def,fun,Prop,Type,variable,where,
    Field,Fintype,Finite,DecidableEq,CommRing,CommSemiring,Ring,CharP,Function,
    Injective,Surjective,Bijective,by,rfl,have,rw,exact,ext,simp,intro,
    iterateFrobenius,RingHom,Commute,convert,apply,using,Nat,Coprime,gcd,
    Odd,Even,IsAPN,IsAB,walsh,Nat,obtain,induction,cases,calc,grind},
  showstringspaces=false,
  columns=fullflexible,
  extendedchars=true,
  frame=single,
  framesep=6pt,
  xleftmargin=4pt,
  aboveskip=12pt,
  belowskip=14pt,
}
\lstset{style=lean}

\newtheorem{lemma}{Lemma}
\newtheorem{theorem}{Theorem}
\newtheorem{corollary}{Corollary}

% boxed "what just happened" note
\newcommand{\note}[1]{%
  \par\smallskip
  \noindent\colorbox{leanbg}{\parbox{\dimexpr\linewidth-2\fboxsep}{\small #1}}%
  \par\smallskip}


\title{\textbf{A divisibility that flattens a spectrum}\\[4pt]
\large Walsh divisibility $\to$ three-valued spectrum $\to$ AB, read slowly}
\author{}
\date{}

\begin{document}
\maketitle

\begin{abstract}
\noindent
One modular fact, on the spectral side.\\[2pt]
$\bullet$\quad \texttt{kasami\_walsh\_div}:\;\;
$2^{(n+1)/2}\mid W(a,b)$ for the Kasami function.\\[6pt]
\noindent
Combined with two sum identities (\textbf{Parseval} and the \textbf{fourth
moment}), the divisibility pins every Walsh value to $\{0,\pm 2^{(n+1)/2}\}$ ---
the definition of \textbf{Almost Bent (AB)}. This is the spectral twin of the
char-2 root lemma: a tiny fact that forces a headline.
\end{abstract}

\bigskip

% =====================================================================
\section*{Setting: the Walsh transform}

\noindent
Let $\chi$ be the sign character of the trace on $F=\mathrm{GF}(2^n)$. The Walsh
coefficient of $f$ at $(a,b)$ is the Fourier coefficient
\[
W(a,b)\;=\;\sum_{x\in F}\chi\bigl(a\,x+b\,f(x)\bigr)\;\in\;\mathbb{Z}.
\]

\begin{lstlisting}
noncomputable def walsh (f : F → F) (a b : F) : ℤ :=
  ∑ x : F, χ (a * x + b * f x)

noncomputable def IsAB {n : ℕ} (_ : Fintype.card F = 2 ^ n) (f : F → F) : Prop :=
  ∀ a : F, a ≠ 0 → ∀ b : F,
    walsh f a b ^ 2 = 0 ∨ walsh f a b ^ 2 = (2 ^ (n + 1) : ℤ)
\end{lstlisting}

\note{\textbf{AB} says the squared spectrum takes only \emph{two} values: $0$ and
$2^{n+1}$. Equivalently $W\in\{0,\pm 2^{(n+1)/2}\}$ --- a three-valued spectrum.}

\bigskip

% =====================================================================
\section{The tiny fact: divisibility}

\begin{lstlisting}
theorem kasami_walsh_div {n : ℕ} (hcard : Fintype.card F = 2 ^ n)
    (k : ℕ) (hk : k ≥ 1) (hcop : Nat.Coprime k n) (hnodd : Odd n) (hn : n ≥ 1)
    (a b : F) :
    (2 : ℤ) ^ ((n + 1) / 2) ∣ walsh (· ^ d k : F → F) a b
\end{lstlisting}

\noindent
\textbf{Why.} Substitute $x=y^{2^k+1}$ (a Gold permutation). The Kasami Walsh sum
becomes a sum over a \emph{quadratic} form,
\[
W(a,b)=\sum_{y}\chi\bigl(a\,y^{2^k+1}+b\,y^{2^{3k}+1}\bigr),
\]
whose third derivative vanishes. A Gauss-sum/McEliece estimate then gives
$2^{(n+1)/2}\mid W$.

\note{This single modular fact is what pins the spectrum to a $2$-adic lattice.
On its own it says nothing about \emph{which} multiples occur --- that needs two
moment identities.}

\bigskip

% =====================================================================
\section{The two moments (short sum identities)}

\noindent
Write the squared spectrum as a histogram. Two sums constrain it.

\subsection*{Parseval (the $L^2$ identity)}
\[
\sum_{b} W(a,b)^2 \;=\; |F|^2 \;=\; (2^n)^2
\qquad (a\neq 0,\;f \text{ bijective}).
\]
\begin{lstlisting}
theorem parseval_perm {n : ℕ} (hcard : Fintype.card F = 2 ^ n)
    (f : F → F) (hf : Function.Bijective f) (a : F) (ha : a ≠ 0) :
    ∑ b : F, walsh f a b ^ 2 = (Fintype.card F : ℤ) ^ 2
\end{lstlisting}

\subsection*{Fourth moment (the $L^4$ identity, from APN)}
\[
\sum_{b} W(a,b)^4 \;=\; 2\,|F|^3 .
\]

\note{Parseval is Plancherel for the additive Fourier transform: orthogonality of
characters ($\sum_x\chi(cx)=0$ for $c\neq 0$) collapses the double sum. The fourth
moment is the autocorrelation/APN identity (the differential ``two-solutions''
shows up as the constant $2$).}

\bigskip

% =====================================================================
\section{The squeeze: lattice $+$ two moments $=$ flat}

\noindent
By divisibility, write $W(a,b)=2^{(n+1)/2}\,k_b$ with $k_b\in\mathbb{Z}$.
Substitute into the two moments (and use $n$ odd):
\[
\sum_b k_b^2 = 2^{\,n-1},
\qquad
\sum_b k_b^4 = 2^{\,n-1}.
\]
Subtract:
\[
\sum_b \bigl(k_b^4-k_b^2\bigr)=\sum_b k_b^2(k_b^2-1)=0 .
\]
Every summand is $\ge 0$, so each is $0$: $k_b^2\in\{0,1\}$. Hence
\[
W(a,b)^2 \in \{0,\;2^{\,n+1}\}.
\]
That is AB. \hfill$\blacksquare$

\begin{lstlisting}
theorem ab_from_moments {n : ℕ} (hcard : Fintype.card F = 2 ^ n)
    (f : F → F) (hodd : Odd n) (hn : n ≥ 1)
    (hparseval : ∀ a, a ≠ 0 → ∑ b, walsh f a b ^ 2 = (Fintype.card F : ℤ) ^ 2)
    (hfourth   : ∀ a, a ≠ 0 → ∑ b, walsh f a b ^ 4 = 2 * (Fintype.card F : ℤ) ^ 3)
    (hdiv      : ∀ a b, (2 : ℤ) ^ ((n + 1) / 2) ∣ walsh f a b) :
    IsAB hcard f
\end{lstlisting}

\note{Integer lattice $+$ two quadratic constraints $\Rightarrow$ the values sit
at $\pm 1$ (after scaling). This is the spectral analogue of ``a product is zero
$\Rightarrow$ a factor is zero'': $k^2(k^2-1)=0$ over a sum of squares.}

\bigskip

% =====================================================================
\section{The headline: APN power permutation, $n$ odd $\Rightarrow$ AB}

\begin{lstlisting}
theorem apn_implies_ab_power {n : ℕ} (hcard : Fintype.card F = 2 ^ n)
    (d : ℕ) (hd : Nat.Coprime d (Fintype.card F - 1))
    (hbij : Function.Bijective (· ^ d : F → F))
    (hodd : Odd n) (hn : n ≥ 1)
    (hapn : IsAPN (· ^ d : F → F)) :
    IsAB hcard (· ^ d : F → F) := by
  apply ab_from_moments hcard _ hodd hn
  · exact fun a ha => parseval_perm hcard _ hbij a ha
  · exact fun a ha => fourth_moment_apn hcard d hbij hapn a ha
  · exact fun a b => walsh_pow_divisibility hcard d hd hodd hn a b
\end{lstlisting}

\note{The structural twin of the differential-side ``transport'' step: the harder
\emph{spectral} result (AB) subsumes the easier \emph{differential} one (APN) when
$n$ is odd. Three small facts --- Parseval, fourth moment, divisibility --- assemble
into the headline.}

\bigskip

% =====================================================================
\section*{The staircase}

\[
\begin{tikzcd}[row sep=large, column sep=large]
\boxed{2^{(n+1)/2}\mid W}
  \arrow[dr, Rightarrow, "{\text{lattice}}"'] & &
\sum W^2{=}|F|^2 \arrow[dl, Rightarrow] \\
& \text{three-valued spectrum} \arrow[d, Rightarrow] &
\sum W^4{=}2|F|^3 \arrow[l, Rightarrow] \\
& \mathrm{IsAB} &
\end{tikzcd}
\]

\bigskip
\subsection*{One-screen summary}

\begin{center}
\renewcommand{\arraystretch}{1.6}
\begin{tabular}{l l l}
\textbf{fact} & \textbf{statement} & \textbf{role} \\
\hline
divisibility & $2^{(n+1)/2}\mid W$ & put $W$ on a lattice \\
Parseval & $\sum_b W^2=|F|^2$ & $\sum k_b^2=2^{n-1}$ \\
4th moment & $\sum_b W^4=2|F|^3$ & $\sum k_b^4=2^{n-1}$ \\
squeeze & $k_b^2(k_b^2-1)=0$ & $W^2\in\{0,2^{n+1}\}$ (AB) \\
\end{tabular}
\end{center}

\bigskip
\noindent
\colorbox{leanbg}{\parbox{\dimexpr\linewidth-2\fboxsep}{\centering
$2^{(n+1)/2}\mid W$ $\;+\;$ Parseval $\;+\;$ 4th moment
$\;\Rightarrow\;$ three-valued spectrum $\;\Rightarrow\;$ AB $\;\Rightarrow\;$ APN.}}

\end{document}
